\section{Problem, scope, and provenance}\label{sec:overview}

For a finite simple graph $G$, let $\doub G$ denote the multigraph in
which every original edge has two parallel copies. A colour class is
\emph{locally irregular} if the degrees of the endpoints of each of its
edges are different in that colour class. Parallel copies may have
different colours. Isolated vertices cause no difficulty. The conjecture
of Grzelec and Wo\'zniak~\cite{GW23} asks whether
\[
 G\text{ connected},\quad G\not\cong K_2
 \quad\Longrightarrow\quad \lir(\doub G)\le2.
\]
All statements below concern finite graphs. Original graphs are simple;
auxiliary connection graphs may have labelled parallel edges and loops.
A loop in a connection graph represents a return path in the original
simple graph, rather than an original loop.

The basic language is an edge weight $x_e\in\{0,1,2\}$ equal to the
number of red copies. The actual colour degrees are
$R(v)=\sum_{e\ni v}x_e$ and $B(v)=2d_G(v)-R(v)$. A comparison is required
in a colour precisely when that colour is present on the edge. Requiring
both comparisons on every edge is a stronger sufficient condition that
some, but not all, of our constructions satisfy.

The original conjecture and results for bipartite graphs appear in
\cite{GW23}; cacti are treated in~\cite{GW24}. The regular-graph result,
a three-colour bound for subcubic graphs, and special subcubic families
are established in~\cite{GOW26}. In its arXiv version, Theorems~14--15
cover independent degree-three vertices and subdivisions of selected
cubic edges into paths of length at least five. We give a full subcubic
proof, retaining its complete finite interfaces. The later sections
address different questions: unbounded-degree constructions, prescribed
boundaries, target feasibility, and limitations of proposed proof methods.

The identifiers A1--A9, B1--B4 and C1--C5 refer to a prior inventory of
research notes; they are preserved for traceability in
Table~\ref{tab:coverage}. They are groups of related assertions, not a
count of mutually independent or historically new theorems. We fold
ordinary existence results for smaller subcubic families into A1 and
retain stronger boundary statements separately. D1--D5 identify completed
extensions documented after 17:00 China Standard Time on 27 October
2026, with their source snapshot fixed at 20:31:04 that evening.
E1--E4 contain four completed notes from 28 October 2026, frozen at
10:25:33 China Standard Time. They add integer rather than merely
modular separation, path/cycle completions, and parity-port certificates.
Work completed after that cutoff, including the later component-dynamic-programming
note, is outside this fixed snapshot.
The original notes, snapshot hashes and a detailed coverage ledger
accompany the manuscript.

\begin{longtable}{>{\raggedright\arraybackslash}p{0.10\linewidth}>{\raggedright\arraybackslash}p{0.53\linewidth}>{\raggedright\arraybackslash}p{0.27\linewidth}}
\caption{Coverage of the retained inventory.}\label{tab:coverage}\\
\toprule Identifier&Assertion group&Location\\\midrule\endfirsthead
\toprule Identifier&Assertion group&Location\\\midrule\endhead
A1&All subcubic graphs; safe interfaces and the two-mark kernel&
Sections~\ref{sec:subcubic}--\ref{sec:complete}, Appendix~\ref{app:cert}\\
A2&Three sparse direct-branch hypotheses&Section~\ref{sec:branches}\\
A3&Degree-homogeneous, three-colourable direct cores&Section~\ref{sec:homogeneous}\\
A4--A5&Balanced half-edges and simultaneous preserved-colour extensions&Section~\ref{sec:balanced}\\
A6&Near-centre parity realisation and degree-gap separation&Section~\ref{sec:parity}\\
A7--A8&Arbitrary clique root sums, block trees, and capacity extensions&Section~\ref{sec:cliques}\\
A9&Tree gluing of given coloured pieces&Section~\ref{sec:pieceglue}\\
B1&Trees, paths, decorated saturation, preserved subdivisions, offsets&Sections~\ref{sec:generalports}, \ref{sec:trees}, \ref{sec:preserve}, \ref{sec:oddforest}\\
B2&Fixed targets, matching, factors, exchanges, defect bounds&Section~\ref{sec:targets}\\
B3&Exact polynomial, free parameters, permanent and parity formula&Section~\ref{sec:algebra}\\
B4&Normal forms, peripheral elimination and minimal counterexamples&Section~\ref{sec:reductions}\\
C1&Cut, parity, modulo-three and modulo-two limitations&Sections~\ref{sec:mod2limits}, \ref{sec:templates}\\
C2&Nonincreasing single-edge and two-unit descent barriers&Section~\ref{sec:descent}\\
C3&Target obstructions, nonconvexity and unbounded parity radius&Sections~\ref{sec:targetbarriers}, \ref{sec:radius}\\
C4&Sharp fixed-interface path and extension thresholds&Sections~\ref{sec:treecapacity}, \ref{sec:cliquecapacity}, \ref{sec:interfacebarriers}\\
C5&Gluing, prescribed root sums and multiple marked ports&Sections~\ref{sec:piecebarrier}, \ref{sec:interfacebarriers}\\
D1&High minimum degree relative to chromatic number&Section~\ref{sec:highdegree}\\
D2&Multiple CRT regions, low-side ports, matched tree attachments&Sections~\ref{sec:crtregions}, \ref{sec:irregularcore}\\
D3&Dense antitone targets: $8\delta\ge7n+6$&Section~\ref{sec:dense}\\
D4&Complement lifts, nonregular families, structural restrictions&Section~\ref{sec:complements}\\
D5&Dense forbidden targets and sequential region absorption&Section~\ref{sec:denselists}\\
E1&Signed gaps with two occurrences; paired partitions and fixed-mixed integer repairs&Section~\ref{sec:signedgap}\\
E2&Degree-three/four path--cycle side against a covering matching&Section~\ref{sec:pathcyclepart}\\
E3&Protected parity joins and exact independent-side completions&Section~\ref{sec:isolatepart}\\
E4&Arbitrary-degree odd-component root ports and capacity-aware parity test&Section~\ref{sec:rootports}\\
\bottomrule
\end{longtable}

The matching reductions, Tutte's factor criterion, jump-system exchange
principles and polynomial interpolation are classical tools, not claims
of invention. We give the required applications and elementary arguments
so that the boundary between a feasible target and a legal colouring is
explicit. The degree-constrained subgraph theorem used for D1 is quoted
with its hypotheses and applied only to a simple host graph. The
subcubic proof uses the degree-choosability theorem recorded
in~\cite{CR17}. Other background dependencies are indicated where used.

This is a research preprint with extensive computer assistance. Its
infinite-family assertions are justified by the displayed arguments and,
where indicated, complete finite kernels plus lifting inductions.
Random and bounded-order checks test implementations; they do not prove
the conjecture. No independent expert review or proof-assistant
verification is claimed, and no absolute historical-priority claim is
made. In particular, none of the strategy obstructions below is a
counterexample to the original conjecture.
